\chapter{Design and Analysis of Single-Factor Experiments: The Analysis of Variance}

\section{Designing Engineering Experiments}

An \emph{experiment} changes controllable inputs on purpose and watches the response.
\begin{itemize}
  \item A \emph{factor} is an input we control. It takes several \emph{levels}.
  \item A \emph{treatment} is one combination of factor levels.
  \item Statistics helps in three places: choosing the design, analysing the data, and stating conclusions with known error rates.
\end{itemize}

Three principles of design:
\begin{itemize}
  \item \textbf{Replication.} Repeat each treatment. This gives an estimate of experimental error and improves precision.
  \item \textbf{Randomisation.} Assign treatments to units, and run the trials, in random order. This protects against hidden systematic effects and justifies the independence assumption. As in Chapter~4, it also supports cause-and-effect conclusions.
  \item \textbf{Blocking.} Group similar units and compare treatments within each group. This removes variation that would otherwise sit in the error (Section~\ref{sec:rcbd}).
\end{itemize}
Before collecting data, state the problem, the response, the factors and ranges, and the planned analysis.

\section{The Completely Randomized Single-Factor Experiment}

\subsection{An Example}

Suppose $a$ treatments (levels of one factor) are compared, and units are assigned to them completely at random. This is a \emph{completely randomized design} (CRD). Let $y_{ij}$ be the $j$-th observation under treatment $i$, $j=1,\dots,n_i$.

\begin{example}\label{ex:anovadata}
Three machine settings (A, B, C) are tried on four randomly assigned runs each. The response is strength.
\begin{center}
\begin{tabular}{lcccc|cc}
\toprule
Setting & \multicolumn{4}{c|}{Observations} & Total $y_{i\cdot}$ & Mean $\bar{y}_{i\cdot}$\\
\midrule
A & 12 & 14 & 13 & 15 & 54 & 13.5\\
B & 16 & 18 & 17 & 19 & 70 & 17.5\\
C & 13 & 15 & 14 & 16 & 58 & 14.5\\
\bottomrule
\end{tabular}
\end{center}
The grand total is $y_{\cdot\cdot}=182$, the grand mean is $\bar{y}_{\cdot\cdot}=15.17$, and $N=12$.
\end{example}

\subsection{The Analysis of Variance}

\begin{definition}[Single-factor model]
\[ Y_{ij}=\mu+\tau_i+\varepsilon_{ij},\qquad i=1,\dots,a;\ j=1,\dots,n_i. \]
\begin{itemize}
  \item $\mu$ is the overall mean and $\tau_i$ is the effect of treatment $i$, with $\sum n_i\tau_i=0$.
  \item The errors $\varepsilon_{ij}$ are independent $\Normal(0,\sigma^2)$.
  \item The treatment means are $\mu_i=\mu+\tau_i$.
\end{itemize}
\end{definition}

The test is
\[ H_0:\mu_1=\mu_2=\dots=\mu_a\ (\text{all }\tau_i=0)\qquad\text{against}\qquad H_1:\mu_i\ne\mu_j\text{ for at least one pair.} \]

\begin{theorem}[Partition of the total variability]\label{thm:anova-ss}
\[ \underbrace{\sum_{i=1}^a\sum_{j=1}^{n_i}(y_{ij}-\bar{y}_{\cdot\cdot})^2}_{SS_T}=\underbrace{\sum_{i=1}^an_i(\bar{y}_{i\cdot}-\bar{y}_{\cdot\cdot})^2}_{SS_{\mathrm{Treatments}}}+\underbrace{\sum_{i=1}^a\sum_{j=1}^{n_i}(y_{ij}-\bar{y}_{i\cdot})^2}_{SS_E}, \]
with degrees of freedom $N-1=(a-1)+(N-a)$.
\end{theorem}

\begin{proof}
Write $y_{ij}-\bar{y}_{\cdot\cdot}=(\bar{y}_{i\cdot}-\bar{y}_{\cdot\cdot})+(y_{ij}-\bar{y}_{i\cdot})$ and square. The cross term is $2\sum_i(\bar{y}_{i\cdot}-\bar{y}_{\cdot\cdot})\sum_j(y_{ij}-\bar{y}_{i\cdot})=0$, because deviations from a group mean sum to zero within each group.
\end{proof}

The \emph{mean squares} are $MS_{\mathrm{Tr}}=SS_{\mathrm{Tr}}/(a-1)$ and $MS_E=SS_E/(N-a)$. Their expected values explain the test:
\[ \Exp[MS_E]=\sigma^2,\qquad\Exp[MS_{\mathrm{Tr}}]=\sigma^2+\frac{\sum n_i\tau_i^2}{a-1}. \]
\begin{itemize}
  \item If $H_0$ is true, both estimate $\sigma^2$, so their ratio is near $1$.
  \item If the means differ, $MS_{\mathrm{Tr}}$ is inflated and the ratio is large.
\end{itemize}

\begin{theorem}[The $F$-test of ANOVA]\label{thm:anova-F}
Under $H_0$ and the model assumptions, $SS_{\mathrm{Tr}}/\sigma^2\sim\chi^2_{a-1}$ and $SS_E/\sigma^2\sim\chi^2_{N-a}$, independently. So
\[ F_0=\frac{MS_{\mathrm{Tr}}}{MS_E}\sim F_{a-1,\,N-a}. \]
Reject $H_0$ if $f_0>f_{\alpha,\,a-1,\,N-a}$ (upper tail only).
\end{theorem}

\begin{center}
\begin{tabular}{lcccc}
\toprule
Source & Sum of squares & df & Mean square & $F_0$ \\
\midrule
Treatments & $SS_{\mathrm{Tr}}$ & $a-1$ & $MS_{\mathrm{Tr}}$ & $MS_{\mathrm{Tr}}/MS_E$ \\
Error & $SS_E$ & $N-a$ & $MS_E$ & \\
Total & $SS_T$ & $N-1$ & & \\
\bottomrule
\end{tabular}
\end{center}

\emph{Computing formulas} (with $y_{i\cdot}$ the treatment totals):
\[ SS_T=\sum\sum y_{ij}^2-\frac{y_{\cdot\cdot}^2}N,\qquad SS_{\mathrm{Tr}}=\sum\frac{y_{i\cdot}^2}{n_i}-\frac{y_{\cdot\cdot}^2}N,\qquad SS_E=SS_T-SS_{\mathrm{Tr}}. \]

\begin{example}[Continuing Example~\ref{ex:anovadata}]\label{ex:anovacalc}
\begin{itemize}
  \item $SS_{\mathrm{Tr}}=4\big[(13.5-15.17)^2+(17.5-15.17)^2+(14.5-15.17)^2\big]=34.67$.
  \item Within each setting the deviations from the mean are $(-1.5,0.5,-0.5,1.5)$. Their squares sum to $5$, so $SS_E=3(5)=15.0$.
  \item $SS_T=49.67$. Degrees of freedom: $2$ and $9$. $MS_{\mathrm{Tr}}=17.33$, $MS_E=1.667$.
  \item $f_0=10.4>f_{0.05,2,9}=4.26$, so reject $H_0$ ($P\approx0.005$). The settings do not all give the same mean strength.
\end{itemize}
\end{example}

The formulas hold for unequal $n_i$. Still, prefer a balanced design ($n_i=n$): the test is more robust to unequal variances, and power is largest for a given $N$.

\begin{example}[Paper tensile strength]
A paper maker tests four hardwood concentrations ($5,10,15,20\%$) with six specimens each, in random order. So $a=4$, $n=6$, $N=24$. The treatment totals are $60,94,102,127$ (means $10.0,\,15.67,\,17.0,\,21.17$) and the grand total is $383$.
\begin{itemize}
  \item $SS_T=512.96$, $SS_{\text{Treatments}}=382.79$, $SS_E=130.17$, with $23$, $3$ and $20$ degrees of freedom.
  \item $MS_{\text{Treatments}}=127.60$, $MS_E=6.51$, $f_0=19.60$.
  \item At $\alpha=0.01$, $f_{0.01,3,20}=4.94$. Reject $H_0$ ($P=3.59\times10^{-6}$).
\end{itemize}
The rejection region is the \emph{upper} tail only, because under $H_1$ the expected value of $MS_{\text{Treatments}}$ exceeds $\sigma^2$.

Intervals: for one treatment mean, $\bar{y}_{i\cdot}\pm t_{\alpha/2,\,a(n-1)}\sqrt{MS_E/n}$. For a difference, $\bar{y}_{i\cdot}-\bar{y}_{j\cdot}\pm t_{\alpha/2,\,a(n-1)}\sqrt{2MS_E/n}$. If the interval for a difference contains $0$, no difference is detected.

Fisher's LSD: $t_{0.025,20}=2.086$, so $\mathrm{LSD}=2.086\sqrt{2(6.51)/6}=3.07$. Every pair of means differs by more than $3.07$, except $10\%$ and $15\%$ (means $15.67$ and $17.0$). So $20\%$ hardwood gives the greatest strength.
\end{example}

\subsection{Multiple Comparisons Following the ANOVA}

A significant $F$ says that \emph{some} means differ. It does not say which.

\begin{itemize}
  \item Comparing all $\binom a2$ pairs with $t$-tests at level $\alpha$ each makes the chance of at least one false positive much larger than $\alpha$. This chance is the \emph{experimentwise error rate}.
  \item For $a=5$ there are $10$ pairs. With $\alpha=0.05$ each, the experimentwise rate is about $0.40$.
\end{itemize}
Two procedures control it.

\begin{definition}[Tukey's method]
In a balanced design, declare $\mu_i\ne\mu_j$ if $\abs{\bar{y}_{i\cdot}-\bar{y}_{j\cdot}}>T_\alpha$, where
\[ T_\alpha=q_\alpha(a,f)\sqrt{\frac{MS_E}n},\qquad f=N-a. \]
$q_\alpha(a,f)$ is the upper $\alpha$ point of the studentised range distribution (tabulated).
\begin{itemize}
  \item Tukey keeps the \emph{overall} error rate for all pairwise comparisons at exactly $\alpha$ (balanced design).
  \item It gives simultaneous intervals $\bar{y}_{i\cdot}-\bar{y}_{j\cdot}\pm T_\alpha$ for all $\mu_i-\mu_j$.
  \item By contrast, the ``family error rate'' that software quotes for unprotected LSD comparisons is the chance of at least one Type~I error among all the comparisons. It is about $0.19$ for $a=4$ (six comparisons).
\end{itemize}
\end{definition}

\begin{definition}[Fisher's least significant difference (LSD)]
Declare $\mu_i\ne\mu_j$ if
\[ \abs{\bar{y}_{i\cdot}-\bar{y}_{j\cdot}}>\mathrm{LSD}=t_{\alpha/2,\,N-a}\sqrt{MS_E\Big(\frac1{n_i}+\frac1{n_j}\Big)}. \]
\begin{itemize}
  \item Apply LSD only \emph{after} a significant $F$-test (``protected LSD''). This protects the error rate.
  \item LSD and Tukey are \emph{two-sided} comparisons, with $\alpha/2$ in each tail. For a one-sided comparison $H_1:\mu_i>\mu_j$, use $t_{\alpha,\,N-a}$ instead of $t_{\alpha/2,\,N-a}$.
  \item The ANOVA $F$-test itself rejects only in the upper tail, $f_0>f_{\alpha,\,a-1,\,N-a}$. The ratio $MS_{\text{Tr}}/MS_E$ becomes large for \emph{any} departure from equal means. There is no ``two-sided $F$-test'' here.
\end{itemize}
\end{definition}

\begin{example}[Continuing Example~\ref{ex:anovacalc}]
\begin{itemize}
  \item Tukey: $q_{0.05}(3,9)=3.95$, so $T_{0.05}=3.95\sqrt{1.667/4}=2.55$.
  \item Differences: $B-A=4.0$, $B-C=3.0$, $C-A=1.0$. Setting B differs from both others. A and C are not distinguished.
  \item LSD: $2.262\sqrt{1.667\cdot2/4}=2.262(0.913)=2.07$. Same conclusions here.
\end{itemize}
\end{example}

\subsection{More About Multiple Comparisons (CD Only)}

Fisher's LSD is easy and very common, but it is liberal. Each test has level $\alpha$, yet the experimentwise Type~I error is much larger than $\alpha$. Three other approaches follow.

\subsubsection*{Graphical comparison of means}

Each treatment average $\bar{y}_{i\cdot}$ has standard deviation $\sigma/\sqrt n$. If all treatment means were equal, the averages would look like a random sample from a normal curve with that spread, estimated by $\sqrt{MS_E/n}$.
\begin{itemize}
  \item Imagine this curve on a movable transparency. Slide it along the axis where the $a$ means are plotted.
  \item If some position makes all the means look like typical draws, there is no evidence of differences.
  \item Means that do not fit under the curve belong to treatments with different mean responses.
\end{itemize}
This is a rough but effective picture of the ANOVA result. (The $t$ distribution would be exact, but the normal curve is close.)

\subsubsection*{Orthogonal contrasts}

\begin{definition}[Contrast]
A \emph{contrast} is a combination of treatment totals $C=\sum_{i=1}^ac_iy_{i\cdot}$ with $\sum c_i=0$ (unbalanced: $\sum n_ic_i=0$). It expresses one specific comparison, chosen \emph{before} the experiment.
\end{definition}
Examples: $H_0:\mu_3=\mu_4$ has coefficients $(0,0,1,-1)$. $H_0:\mu_1+\mu_3=\mu_2+\mu_4$ has coefficients $(1,-1,1,-1)$.

The sum of squares of a contrast has one degree of freedom:
\[ SS_C=\frac{\big(\sum_ic_iy_{i\cdot}\big)^2}{n\sum_ic_i^2}\qquad\Big(\text{unbalanced: }\frac{\big(\sum c_iy_{i\cdot}\big)^2}{\sum n_ic_i^2}\Big). \]
It is tested by $F_0=SS_C/MS_E\sim F_{1,\,N-a}$.

\begin{itemize}
  \item Contrasts $\{c_i\}$ and $\{d_i\}$ are \emph{orthogonal} if $\sum c_id_i=0$ (unbalanced: $\sum n_ic_id_i=0$).
  \item A set of $a-1$ mutually orthogonal contrasts splits $SS_{\mathrm{Tr}}$ into $a-1$ independent one-degree-of-freedom pieces. The tests are independent.
  \item With $a=3$ and treatment 1 a control, $(-2,1,1)$ (control against the average of the others) and $(0,1,-1)$ (one level against the other) are orthogonal.
  \item Choose contrasts \emph{before} looking at the data. Picking the largest observed differences afterwards inflates the Type~I error.
\end{itemize}

\begin{example}
For Example~\ref{ex:anovacalc}: $y_{1\cdot}=54$, $y_{2\cdot}=70$, $y_{3\cdot}=58$, $n=4$.
\begin{itemize}
  \item $c=(-2,1,1)$: $C=-108+70+58=20$ and $SS_C=20^2/(4\cdot6)=16.67$.
  \item $d=(0,1,-1)$: $C=12$ and $SS_d=12^2/(4\cdot2)=18.00$.
  \item $c\cdot d=0$, so they are orthogonal, and $SS_c+SS_d=34.67=SS_{\mathrm{Tr}}$.
  \item Against $MS_E=1.667$: $F_c=10.0$ and $F_d=10.8$. Compare each with $f_{0.05,1,9}=5.12$. Both are significant.
\end{itemize}
\end{example}

\subsubsection*{Tukey's method}

Tukey's procedure compares all pairs using the studentised range
\[ Q=\frac{\bar{Y}_{\max}-\bar{Y}_{\min}}{\sqrt{MS_E/n}}, \]
where $\bar{Y}_{\max}$ and $\bar{Y}_{\min}$ are the largest and smallest of the $a$ sample means. For equal sample sizes, declare $\mu_i\ne\mu_j$ if $\abs{\bar{y}_{i\cdot}-\bar{y}_{j\cdot}}>T_\alpha=q_\alpha(a,f)\sqrt{MS_E/n}$, with $f=N-a$. The overall error rate for all pairwise comparisons is exactly $\alpha$.

\begin{interview}{beyond the textbook}
Other procedures you will meet: Bonferroni (test each of $m$ comparisons at $\alpha/m$), Scheff\'e (all possible contrasts; most conservative), Dunnett (each treatment against a control), and the Holm and Benjamini--Hochberg corrections (Chapter~13). The textbook covers Fisher's LSD, the graphical method, orthogonal contrasts and Tukey's method.
\end{interview}

\subsection{Residual Analysis and Model Checking}

The residuals are $e_{ij}=y_{ij}-\bar{y}_{i\cdot}$. Check the model as in regression.
\begin{itemize}
  \item \textbf{Normal probability plot of residuals.} Checks normality. The $F$-test is fairly robust to moderate departures.
  \item \textbf{Residuals against fitted values $\bar{y}_{i\cdot}$.} A funnel shape means the variance grows with the mean. Try $\ln y$ or $\sqrt y$.
  \item \textbf{Residuals against run order.} A trend suggests a time effect, or that randomisation or independence failed.
  \item \textbf{Equal variances (beyond the textbook).} \emph{Bartlett's test} of $H_0:\sigma_1^2=\dots=\sigma_a^2$ uses $\chi_0^2=2.3026\,q/c$, with
  \[ q=(N-a)\log_{10}S_p^2-\sum(n_i-1)\log_{10}S_i^2,\qquad c=1+\frac1{3(a-1)}\Big(\sum\frac1{n_i-1}-\frac1{N-a}\Big). \]
  Compare with $\chi^2_{\alpha,\,a-1}$. It is sensitive to non-normality, so Levene's test is often preferred.
\end{itemize}

\subsection{Determining Sample Size}

The power of the $F$-test depends on a noncentrality parameter. For the balanced case define
\[ \Phi^2=\frac{n\sum_{i=1}^a\tau_i^2}{a\sigma^2}. \]
Read the Type~II error $\beta$ from OC curves with $a-1$ and $a(n-1)$ degrees of freedom. (Or compute power from the noncentral $F$ with noncentrality $\lambda=a\Phi^2$.)

To plan the experiment:
\begin{enumerate}
  \item Choose $D$, the smallest range of treatment means (largest minus smallest) that matters in practice.
  \item The worst case for a given $D$ puts two means at the extremes and the rest at the centre. Then $\Phi^2=nD^2/(2a\sigma^2)$.
  \item Use a prior estimate of $\sigma$. Choose $n$ so the power is adequate (say $0.8$ or $0.9$).
\end{enumerate}
Larger $n$, larger $D$ and smaller $\sigma$ all increase power.

\begin{example}
Five means, $\alpha=0.01$. We want to reject $H_0$ with probability at least $0.90$ when $\sum\tau_i^2=5.0\sigma^2$.
\begin{itemize}
  \item $\Phi^2=n\sum\tau_i^2/(a\sigma^2)=n$, with $\nu_1=4$ and $\nu_2=5(n-1)$.
  \item Try $n=4$: $\Phi^2=4$, $\Phi=2$, $\nu_2=15$. The OC curve gives $\beta\approx0.38$, so power $\approx0.62<0.90$. Not enough.
  \item Continue the same way. At least $n=6$ replicates are needed.
\end{itemize}
\end{example}

\subsection{Technical Details about the Analysis of Variance (CD Only)}

Write the observations as a vector. Each sum of squares is a quadratic form $\mathbf{y}\T A_k\mathbf{y}$, where the $A_k$ are symmetric idempotent matrices with $\sum A_k=I$ and $A_jA_k=0$.

By Cochran's theorem, the forms $\mathbf{y}\T A_k\mathbf{y}/\sigma^2$ are then independent noncentral chi-square variables. Each has degrees of freedom $\operatorname{rank}A_k$ and noncentrality $\boldsymbol\mu\T A_k\boldsymbol\mu/\sigma^2$. The noncentrality is zero exactly when that effect is absent.

Geometrically, the observation vector is split into orthogonal parts: the grand mean, the treatment effect and the residual (Linear Algebra, Chapters~11--12). The ratio of two independent scaled central chi-squares is $F$. This proves Theorem~\ref{thm:anova-F}.

\section{The Random Effects Model}

\subsection{Fixed Versus Random Factors}

\begin{itemize}
  \item \textbf{Fixed effects.} The $a$ treatments are the specific levels of interest. Conclusions refer to those levels only.
  \item \textbf{Random effects.} The $a$ levels are a random sample from a large population of levels (for example, $a$ operators, batches or looms chosen at random). We want to describe the whole population, in particular how much variability the factor contributes.
\end{itemize}

\subsection{ANOVA and Variance Components}

\begin{definition}[Random effects model]
$Y_{ij}=\mu+\tau_i+\varepsilon_{ij}$ with $\tau_i\sim\Normal(0,\sigma_\tau^2)$ and $\varepsilon_{ij}\sim\Normal(0,\sigma^2)$, all independent.
\begin{itemize}
  \item The \emph{variance components} are $\sigma_\tau^2$ and $\sigma^2$.
  \item $\Var(Y_{ij})=\sigma_\tau^2+\sigma^2$.
\end{itemize}
\end{definition}

The ANOVA table and sums of squares are exactly as before. The hypothesis and the interpretation change.
\begin{itemize}
  \item Test $H_0:\sigma_\tau^2=0$ against $H_1:\sigma_\tau^2>0$ with the same $F_0=MS_{\mathrm{Tr}}/MS_E\sim F_{a-1,\,N-a}$ under $H_0$.
  \item Testing individual effects $\tau_i=0$ is meaningless here. Also $\sum\tau_i=0$ does not apply.
\end{itemize}
For balanced data ($n_i=n$) the expected mean squares are
\[ \Exp[MS_E]=\sigma^2,\qquad\Exp[MS_{\mathrm{Tr}}]=\sigma^2+n\sigma_\tau^2. \]
So the method-of-moments (\emph{ANOVA}) estimators are
\[ \hat\sigma^2=MS_E,\qquad\hat\sigma_\tau^2=\frac{MS_{\mathrm{Tr}}-MS_E}n. \]

\begin{example}
Treat the settings of Example~\ref{ex:anovacalc} as a random sample of levels.
\begin{itemize}
  \item $\hat\sigma^2=1.667$ and $\hat\sigma_\tau^2=(17.33-1.667)/4=3.92$.
  \item The share of total variance due to the factor (the \emph{intraclass correlation}) is $3.92/(3.92+1.667)=0.70$.
\end{itemize}
\end{example}

\begin{example}[Looms]
A textile company samples four looms at random and makes four strength determinations per loom. The ANOVA shows significant loom-to-loom variability.
\begin{itemize}
  \item $\hat\sigma^2=1.90$ (within looms) and $\hat\sigma_\tau^2=6.96$ (between looms).
  \item $\widehat{V}(Y_{ij})=6.96+1.90=8.86$. The process standard deviation is $\sqrt{8.86}=2.98$ psi about a mean of $95.45$ psi. Most of the variability comes from differences between looms.
  \item If the loom-to-loom variability were removed, the standard deviation would fall to $\sqrt{1.90}=1.38$ psi. That would greatly reduce fallout below a lower specification limit of $90$ psi.
\end{itemize}
\end{example}

\begin{pitfall}{negative variance estimates}
If $MS_{\mathrm{Tr}}<MS_E$, the ANOVA estimate $\hat\sigma_\tau^2$ is negative. This suggests the true $\sigma_\tau^2$ is near or at zero. The textbook lists three courses of action:
\begin{itemize}
  \item Accept it as evidence that the true component is zero (this disturbs the properties of the other estimates). In practice, report $\hat\sigma_\tau^2=0$.
  \item Re-estimate with a method that always gives nonnegative values, such as maximum likelihood or REML.
  \item Treat it as a sign that the assumed linear model is wrong, and study the model.
\end{itemize}
\end{pitfall}

\subsection{Determining Sample Size in the Random Model (CD Only)}

Under $H_1$, the statistic $F_0$ has a \emph{central} $F$ distribution, scaled by $\lambda=\sqrt{1+n\sigma_\tau^2/\sigma^2}$:
\[ F_0/\lambda^2\sim F_{a-1,\,a(n-1)}. \]
So the power is $\Prob\big(F_{a-1,\,a(n-1)}>f_{\alpha}/\lambda^2\big)$. It grows with $n\sigma_\tau^2/\sigma^2$.

The levels are themselves sampled. So at a fixed total $N=an$, adding levels $a$ usually helps more than adding replication $n$.

\section{Randomized Complete Block Design}\label{sec:rcbd}

\subsection{Design and Statistical Analysis}

When units are different (different days, batches, operators), treatment comparisons in a CRD get mixed up with that variation. This inflates $MS_E$.

In a \emph{randomized complete block design} (RCBD):
\begin{itemize}
  \item Units are grouped into $b$ \emph{blocks} of $a$ similar units.
  \item Within each block, each of the $a$ treatments appears \emph{once}, in random order.
  \item Treatments are compared within blocks, where the units are alike.
\end{itemize}

\begin{definition}[RCBD model]
$Y_{ij}=\mu+\tau_i+\beta_j+\varepsilon_{ij}$, with $i=1,\dots,a$ (treatments) and $j=1,\dots,b$ (blocks). Also $\sum\tau_i=0$, $\sum\beta_j=0$, and $\varepsilon_{ij}\sim\Normal(0,\sigma^2)$ independent.
\end{definition}

\begin{theorem}[Partition for the RCBD]\label{thm:rcbd-ss}
\[ SS_T=SS_{\mathrm{Tr}}+SS_{\mathrm{Blocks}}+SS_E, \]
with
\begin{align*}
SS_{\mathrm{Tr}}&=b\sum_i(\bar{y}_{i\cdot}-\bar{y}_{\cdot\cdot})^2,\\
SS_{\mathrm{Blocks}}&=a\sum_j(\bar{y}_{\cdot j}-\bar{y}_{\cdot\cdot})^2,\\
SS_E&=\sum_{i,j}(y_{ij}-\bar{y}_{i\cdot}-\bar{y}_{\cdot j}+\bar{y}_{\cdot\cdot})^2,
\end{align*}
and degrees of freedom $ab-1=(a-1)+(b-1)+(a-1)(b-1)$.
\end{theorem}

\begin{theorem}[Tests for the RCBD]\label{thm:rcbd-F}
To test $H_0:\tau_1=\dots=\tau_a=0$, use
\[ F_0=\frac{MS_{\mathrm{Tr}}}{MS_E}\sim F_{a-1,\,(a-1)(b-1)}\ \text{ under }H_0. \]
Reject if $f_0>f_{\alpha,\,a-1,\,(a-1)(b-1)}$.

Blocks can be tested by $MS_{\mathrm{Blocks}}/MS_E$. This is of secondary interest: blocks were not randomised, so it is only a rough guide to whether blocking was worthwhile.
\end{theorem}

Computing formulas, with totals $y_{i\cdot}$ (treatments), $y_{\cdot j}$ (blocks) and $y_{\cdot\cdot}$:
\[ SS_{\mathrm{Tr}}=\frac1b\sum y_{i\cdot}^2-\frac{y_{\cdot\cdot}^2}{ab},\qquad SS_{\mathrm{Blocks}}=\frac1a\sum y_{\cdot j}^2-\frac{y_{\cdot\cdot}^2}{ab}. \]

\begin{example}\label{ex:rcbd}
Three treatments in four blocks (rows are treatments, columns are blocks):
\begin{center}
\begin{tabular}{l cccc|c}
\toprule
 & B1 & B2 & B3 & B4 & $y_{i\cdot}$\\
\midrule
T1 & 10 & 12 & 11 & 15 & 48\\
T2 & 13 & 14 & 13 & 18 & 58\\
T3 & 12 & 12 & 13 & 16 & 53\\
\midrule
$y_{\cdot j}$ & 35 & 38 & 37 & 49 & 159\\
\bottomrule
\end{tabular}
\end{center}
\begin{itemize}
  \item Correction term: $159^2/12=2106.75$. With $\sum y^2=2161$, $SS_T=54.25$.
  \item $SS_{\mathrm{Tr}}=(48^2+58^2+53^2)/4-2106.75=12.50$.
  \item $SS_{\mathrm{Blocks}}=(35^2+38^2+37^2+49^2)/3-2106.75=39.58$.
  \item $SS_E=54.25-12.50-39.58=2.17$.
  \item Degrees of freedom $2,3,6$: $MS_{\mathrm{Tr}}=6.25$, $MS_{\mathrm{Blocks}}=13.19$, $MS_E=0.361$.
  \item $f_0=17.3>f_{0.05,2,6}=5.14$, so the treatment means differ. The block ratio $F=36.5$ shows large block-to-block differences.
\end{itemize}
\end{example}

\subsection{Multiple Comparisons}

Each treatment mean is an average of $b$ observations. So:
\[ \text{Tukey: }T_\alpha=q_\alpha\big(a,(a-1)(b-1)\big)\sqrt{MS_E/b},\qquad\text{LSD: }t_{\alpha/2,\,(a-1)(b-1)}\sqrt{2MS_E/b}. \]

\begin{example}
For Example~\ref{ex:rcbd}: $q_{0.05}(3,6)=4.34$ and $T=4.34\sqrt{0.361/4}=1.30$.

The treatment means are $12.0$, $14.5$, $13.25$. T2 differs from T1 (difference $2.5$). T3 is not distinguished from either (differences $1.25$ and $1.25$).
\end{example}

\begin{example}[Fabric strength]
Four chemicals ($a=4$) are tested once each on five fabric samples ($b=5$, the blocks), in random order within each sample. The error has $(a-1)(b-1)=12$ degrees of freedom.
\begin{itemize}
  \item $f_0=75.13>f_{0.05,3,12}=5.95$, so the chemicals differ.
  \item Fisher's LSD (each average uses $b=5$ observations) is $0.39$. Any pair of averages differing by $0.39$ or more is significantly different.
  \item Chemical~4 differs from the other three. Chemicals $2$ and $3$ do not differ; neither do $1$ and $3$.
  \item Without blocking, the error would have $16$ degrees of freedom, so blocking cost only $b-1=4$. The block effect was large and would otherwise have inflated $MS_E$.
\end{itemize}
Rule of thumb: when in doubt whether block effects matter, block. The block $F$-ratio that software prints is of doubtful validity as a formal test, because randomisation was restricted to within blocks.
\end{example}

\begin{interview}{was blocking worth it?}
Compare with a CRD on the same units. The error mean square a CRD would have had is estimated by
\[ \widehat{MS}_E^{\mathrm{CRD}}=\frac{(b-1)MS_{\mathrm{Blocks}}+b(a-1)MS_E}{ab-1}. \]
The relative efficiency is $\widehat{MS}_E^{\mathrm{CRD}}/MS_E$.

In Example~\ref{ex:rcbd}: $\big(3(13.19)+4\cdot2(0.361)\big)/11=3.86$, and $3.86/0.361=10.7$. A CRD would need about ten times as many units for the same precision.

Blocking costs $b-1$ error degrees of freedom. It can hurt only if the blocks do not actually reduce variability.
\end{interview}

\subsection{Residual Analysis and Model Checking}

The residuals are $e_{ij}=y_{ij}-\bar{y}_{i\cdot}-\bar{y}_{\cdot j}+\bar{y}_{\cdot\cdot}$. Use the usual plots: normal probability, residuals against fitted values, against treatment, against block and against time.

The RCBD assumes \emph{no treatment-by-block interaction}: treatment effects are the same in every block.
\begin{itemize}
  \item Non-additivity shows up as a pattern in the residual plots. Tukey's one-degree-of-freedom test for nonadditivity can detect it.
  \item If the interaction is real, rethink the model, for example with replication within blocks (Chapter~8).
\end{itemize}

\subsection{Randomized Complete Block Design with Random Factors (CD Only)}

If the blocks (or the treatments) are a random sample, the model has random effects $\beta_j\sim\Normal(0,\sigma_\beta^2)$. The ANOVA table and the treatment $F$-test do not change. With fixed treatments the expected mean squares are
\[ \Exp[MS_E]=\sigma^2,\qquad\Exp[MS_{\mathrm{Blocks}}]=\sigma^2+a\sigma_\beta^2,\qquad\Exp[MS_{\mathrm{Tr}}]=\sigma^2+\frac{b\sum\tau_i^2}{a-1}. \]
So the block variance component is $\hat\sigma_\beta^2=(MS_{\mathrm{Blocks}}-MS_E)/a$.

For Example~\ref{ex:rcbd}: $\hat\sigma_\beta^2=(13.19-0.361)/3=4.28$ and $\hat\sigma^2=0.361$. Nearly all the variability is between blocks.

\begin{summarybox}
\begin{itemize}
  \item Basic design principles: replication, randomisation, blocking. CRD model $y_{ij}=\mu+\tau_i+\varepsilon_{ij}$; $SS_T=SS_{\mathrm{Tr}}+SS_E$; $F_0=MS_{\mathrm{Tr}}/MS_E\sim F_{a-1,N-a}$ (upper tail).
  \item After a significant $F$, use Tukey ($T_\alpha=q_\alpha(a,N-a)\sqrt{MS_E/n}$), Fisher's LSD, Bonferroni, Dunnett or Scheff\'e. Many separate $t$-tests inflate the error rate.
  \item Check normality, constant variance (Bartlett, Levene) and independence (run order) from residuals. Transform the response if needed.
  \item Random effects: same $F$-test, now of $H_0:\sigma_\tau^2=0$. $\hat\sigma_\tau^2=(MS_{\mathrm{Tr}}-MS_E)/n$.
  \item RCBD: $SS_T=SS_{\mathrm{Tr}}+SS_{\mathrm{Blocks}}+SS_E$ with $(a-1)(b-1)$ error df. Compare treatments within blocks. Efficiency against a CRD comes from $\widehat{MS}_E^{\mathrm{CRD}}$.
\end{itemize}
\end{summarybox}

\section{Exercises}

\begin{exercise}
Four methods of teaching are compared, with scores: method 1: $72,75,78,70,74$; method 2: $80,84,79,85,82$; method 3: $74,76,73,78,75$; method 4: $68,71,70,66,69$. Construct the ANOVA table, test equality of the means at $\alpha=0.05$, and use Tukey's method to find which methods differ ($q_{0.05}(4,16)=4.05$).
\end{exercise}

\begin{exercise}
For Exercise~1, compute the residuals, make a normal probability plot, and test equality of variances with Bartlett's test.
\end{exercise}

\begin{exercise}
Derive $\Exp[MS_E]=\sigma^2$ for the single-factor model by showing that $SS_E/\sigma^2$ is a sum of $a$ independent $\chi^2_{n_i-1}$ variables.
\end{exercise}

\begin{exercise}
Five operators are randomly selected from a large pool and each produces four items; $MS_{\mathrm{Operators}}=48$ and $MS_E=12$. Test $H_0:\sigma_\tau^2=0$ ($f_{0.05,4,15}=3.06$) and estimate $\sigma_\tau^2$, $\sigma^2$ and the proportion of variance due to operators.
\end{exercise}

\begin{exercise}
Four fertilisers are tested on plots in five fields (blocks); the ANOVA gives $SS_{\mathrm{Tr}}=120$, $SS_{\mathrm{Blocks}}=300$, $SS_E=90$. Complete the table, test treatments at $\alpha=0.05$, and estimate the relative efficiency of the RCBD to a CRD.
\end{exercise}

\begin{exercise}
Show that in a balanced design with $a=2$ treatments the ANOVA $F$-test is equivalent to the pooled two-sample $t$-test ($F_0=T_0^2$).
\end{exercise}

\begin{exercise}
An experimenter runs four treatments on $24$ units arranged in a CRD but discovers the units fall into six identifiable batches. Explain how the analysis should change, and why ignoring the batches affects the $F$-test.
\end{exercise}

\begin{exercise}
For $a=4$ treatments and $\sigma=2$, find (using $\Phi^2=nD^2/(2a\sigma^2)$) the replication $n$ needed to detect a range of means $D=4$ with high power; discuss how the answer would change with $\sigma=3$.
\end{exercise}
